\documentclass[a4paper]{article}

% Options possibles : 10pt, 11pt, 12pt (taille de la fonte)
%                     oneside, twoside (recto simple, recto-verso)
%                     draft, final (stade de développement)

\usepackage[utf8]{inputenc}   % LaTeX, comprends les accents !
\usepackage[T1]{fontenc}      % Police contenant les caractères français
\usepackage[francais]{babel}
\usepackage{fullpage}
\usepackage{graphicx}  % pour inclure des images
\usepackage{xcolor}    % pour utiliser des couleurs

%\pagestyle{headings}        % Pour mettre des entêtes avec les titres
                                % des sections en haut de page

\title{Projet de programmation L3\\[0.5cm]
  Développement de méthodes pour un logiciel d'extraction de courbes temps-activité vasculaire en tomographie d'émission de positons (TEP)}
\author{Alia CHEBBI, Camille GUIGUEN, DIEVART Colin, Imane LAZIZI, Kanzy YOUSSEF\\[0.5cm]
    \emph{Encadrant :} M. Kucharczak\\
    \emph{En lien avec :} CHU de Montpellier\\[0.5cm]
    \emph{L3 informatique}\\
    Faculté des Sciences\\
    Université de Montpellier.
    }
\date{2024-2025}

\begin{document}

\maketitle                    % Faire un titre utilisant les données
                                % passées à \title, \author et \date

%\begin{center}               % pour centrer
% \includegraphics[scale=1]{}   % insertion d'une image
%\end{center}

\begin{abstract}     % Résumé du travail
    \textcolor{blue}{\emph{Description très succincte du problème et des différentes étapes de réalisation}}
\end{abstract}
\newpage
%\dominitoc  % initializer les minitoc
\tableofcontents
\newpage


\section*{Introduction} % Commencer une section, * permet de ne pas numéroter l'introduction
    \emph{1-2 pages}
    %Présentation du Sujet (1-2 pages)
    \textcolor{blue}{
    \begin{itemize}% pour faire une liste à puces
        \item Présenter le problème étudié et le contexte dans lequel le projet se positionne.
        \item Motiver l'intérêt du problème étudié par rapport à votre parcours d'études et au monde de l'informatique.
        \item Présenter les différentes approches possibles pour la résolution du problème, et en particulier celle choisie.
        \item Donner le cahier des charges détaillé
    \end{itemize}
    }
\emph{Ajout Camille version 1 intro}\\[0.5cm]
    Dans le cadre de ce projet de programmation, nous avons collaboré avec le service de médecine nucléaire 
    du CHU de Montpellier pour implémenter des méthodes à ajouter à un logiciel.\\

    Ce logiciel vise à traiter des données dynamiques obtenues par Tomographie par Émission de Positons (TEP), 
    une technique d'imagerie médicale permettant de visualiser et de mesurer l'activité métabolique des organes 
    et des tissus du corps. 
    On injecte au patient un médicament contenant des traceurs radioactifs envoyant des signaux pour créer des images 
    tridimensionnelles représentant la distribution du radiotraceur dans les tissus.
    Les images reconstruites sont alors enregistrées sous forme de fichiers au format NifTI.\\

    Il a été demandé à notre groupe d’implémenter des méthodes automatisées ou semi-automatiques à ajouter au logiciel 
    de M. Kucharczak. Ces méthodes doivent estimer des Courbes d'entrée vasculaire Temps-Activité (TAC) dans la région de 
    la carotide. Entre autres, analyser la concentration du médicament dans la carotide grâce aux données des fichiers 
    Nifti qui nous ont été fournis.\\

    Ce projet en particulier a attiré notre attention car il se distingue des projets de programmation standards comme le 
    développement de jeux. Il nous permet de découvrir l’application de l’informatique dans un domaine nouveau qui nous 
    intéresse, la médecine nucléaire. De plus, cela nous donne un bon exemple de l’utilité parfois vitale de la programmation.\\

    \vspace{0.5cm}
\noindent\underline{\textbf{CAHIER DES CHARGES :}}

    \vspace{0.2cm}
    \noindent\textbf{Objectif final :}\\
    Implémentation de méthode(s) provenant d’un algorithme que nous jugeons optimal par rapport aux pistes que nous avons étudiées.
    \begin{itemize}
        \item \textbf{En entrée :} un fichier au format NifTI contenant des données patients déjà recadrées de la zone du cerveau et comprenant la zone carotidienne.
        \item \textbf{En sortie :} une ou plusieurs courbes temps-activité (TAC) caractéristiques de la zone carotidienne.
    \end{itemize}

    \vspace{0.2cm}
    \noindent\textbf{Problématiques :}
    \begin{itemize}
        \item Comment extraire la zone carotidienne des images TEP
        \item Les patients "anormaux"
    \end{itemize}

\newpage
\textcolor{blue}{\section{Baert plan de base}}
    \emph{je laisse pour garder une idee de ce quon doit dire}
    \subsection{Développements Logiciel : Conception, Modélisation, Implémentation}
    \emph{8 à 15 pages}
    \textcolor{blue}{
    Il se peut que la conception du logiciel et l'analyse algorithmique aient un poids différent dans votre projet. Les points décrits ci-dessous seront ainsi plus ou moins développés en fonction du travail réalisé en accord avec l'encadrant.
    \begin{enumerate}% pour faire une liste numérotée
        \item Présenter les développements logiciel réalisés dans le cadre du projet.
        \item Présenter les principaux modules du logiciel développé dans le cadre du projet. Utiliser le langage UML pour la modélisation : donner le diagramme de cas d'utilisation et le diagramme des classes.
        \item Décrire les fonctionnalités de l'interface graphique implémentée (si votre logiciel dispose d'une interface graphique).
        \item Décrire le format des données en entrée ou encore les conventions utilisées pour les entrées de vos programmes. Décrire les procédures de lecture et validation des entrées.
        \item Statistiques : nombre de modules/composantes/classes/scripts développés. Nombre de lignes de code.
    \end{enumerate}
    }
    \subsection{Algorithmes et Structures de Données}
    \emph{2 à 3 pages}
    \textcolor{blue}{
    \begin{enumerate}% pour faire une liste numérotée
        \item Présenter les principales structures de données définies dans le cadre du projet.
        \item Présenter les principaux algorithmes implémentés. En illustrer le fonctionnement avec des exemples. Attention : il s'agit ici de choisir 1 ou 2 algorithmes intéressants, et non pas de présenter tous les algorithmes implémentés.
        \item Évaluer la complexité théorique en temps des algorithmes présentés
    \end{enumerate}
    }

\newpage
\section{Algorithmes et chemins de pensée}
    \subsection{Trouver une méthode générale}
        \begin{itemize}
            \item article sprint 1:
            \emph{Ajout Alia version 1}\\
            
            \noindent
            Un des plus grands défis de ce projet a été, pour les informaticiens en herbe que nous sommes, de comprendre les objectifs du projet d’un point de vue scientifique. Nous avons pris appui sur les articles de référence que notre encadrant nous a transmis (voir l'article de Volpi et al., 2023), et nous avons essayé de nous familiariser avec le logiciel \textbf{ITK-SNAP} (voir la section ``blabla'' pour plus de détails).
            
            \noindent
            Afin d’aborder la problématique des images TEP dynamiques en 4D, nous avons d’abord exploré manuellement les fichiers d’acquisition. Grâce à une première série de scripts Python, nous avons visualisé les images en coupe axiale, coronale et sagittale du fichier \texttt{w\_PET\_4D.nii} du patient P01. Pour mieux appréhender la dimension temporelle, nous avons extrait les valeurs d’intensité pour chaque voxel $(x, y, z, t)$, trié les intensités supérieures à un seuil, et sauvegardé ces données volumineuses dans un fichier brut (près de 7 millions de lignes). Ce fichier nous a servi de base de comparaison pour valider les courbes temps-activité (TAC) obtenues ensuite.
            
            \noindent
            Après cette phase de documentation, et la lecture attentive des articles comme \emph{An update on the use of image-derived input functions for human PET studies: new hopes or old illusions?}, plusieurs méthodes d’estimation de la fonction d’entrée nous ont été présentées :

            \begin{enumerate}
                \item \textbf{SIME (Simultaneous Estimation)} : une méthode qui estime simultanément les paramètres cinétiques et la fonction d’entrée artérielle à partir des courbes temps-activité (TAC) issues de l’image ;
                \item \textbf{PBIF (Population-Based Input Function)} : une approche qui, au lieu de mesurer directement l'activité du traceur dans le sang, utilise une courbe moyenne dérivée d’une population de référence, adaptée au sujet par des échantillons de sang ponctuels ;
                \item \textbf{IDIF (Image-Derived Input Function)} : une méthode non invasive basée sur la segmentation de structures vasculaires dans l’image, et le \textit{clustering} des données dynamiques. C’est cette dernière que nous avons choisie d’explorer et d’implémenter comme méthode principale.
            \end{enumerate}
            
            Parmi les différentes variantes testées de l’IDIF, nous avons estimé que l’approche basée sur le \textit{clustering} était la plus adaptée à notre cas. Cet algorithme consiste à regrouper les zones d’intérêt (ROIs) de l’image TEP en fonction de similarités temporelles et spatiales, selon des critères prédéfinis.
            
            Le processus d'extraction de l'IDIF via \textit{clustering} se déroule en trois grandes étapes :
            
            \begin{enumerate}
                \item \textbf{Clustering} des voxels pour regrouper les régions dynamiquement similaires ;
                \item \textbf{Segmentation} des vaisseaux d’intérêt à partir de ces clusters ;
                \item \textbf{Correction}, notamment des effets de volume partiel (PVE), pour obtenir une estimation plus fidèle de la courbe d’entrée.
            \end{enumerate}
            
            Nous avons également comparé différentes approches d’affichage des TAC, en analysant les courbes extraites pour un voxel spécifique et pour un cube de voxels avoisinants, afin d’évaluer la stabilité et la représentativité du signal. En parallèle, nous avons développé une méthode d’extraction semi-automatique des valeurs de la carotide, d’abord sans bibliothèque externe (dans le fichier \texttt{test.py}), puis en intégrant des outils comme \textbf{Nilearn} pour l’automatisation de la segmentation, l’extraction de masques, et le traitement spatio-temporel. Nous avons également exploré le module de parcellisation de Nilearn, ouvrant des perspectives intéressantes pour une segmentation plus robuste et intelligente. \\

            \item imane piste without clustering
            \item choix logique du clustering
        \end{itemize}

    \subsection{les différents clustering}
        \begin{itemize}
            \item Kmeans
            \item Means shift
            \item ward (rena)
        \end{itemize}

    \subsection{MergedMethod}

\newpage
\section{Outils et Environnement de Développement}
    \emph{Ajout Camille version 1}\\
    Pour mener à bien notre projet, nous avons utilisé différents outils.\\

    \noindent
    Notre source de données est constituée de fichiers NIfTI contenant les données d’imagerie TEP des différents patients. 
    Pour visualiser ces fichiers, nous avons utilisé le logiciel ITK-SNAP, conseillé par notre encadrant. 
    ITK-SNAP est un logiciel open-source spécialisé dans la segmentation d'images médicales. Il offre une visualisation des données sur 
    des coupes axiales, transversales et sagitales du cerveau ainsi que les differentes frames ce qui nous permet d'observer l'évolution de l'actif.
    \textcolor{pink}{(jsp quoi dire sur itksnap)}\\

    \noindent
    Nous avons codé en Python pour ce projet. Le logiciel étant déjà en Python, il était obligatoire de fournir des méthodes dans 
    le même langage. Python est idéal pour ce type de projet. C’est un langage très flexible et simple d'accès pour nous ou des non 
    programmeurs voulant reprendre le code. Le langage possède également un grand nombre de bibliothèques spécialisées qui nous ont été essentielles.\\
    \noindent

    \textcolor{pink}{je connais pas bien les bibliothèques. à voir avec vous pour + de détails. la c'est juste chat generated}\\
    Parmi les bibliothèques utilisées, Nilearn nous a permis de manipuler et de visualiser les fichiers NIfTI, facilitant ainsi l'extraction et 
    l'analyse des données TEP. 
    Scipy.ndimage a été essentielle pour appliquer divers filtres et transformations aux données d'imagerie, 
    tandis que Skimage nous a fourni des outils puissants pour la segmentation et la correction de mouvement. Ces bibliothèques ont grandement contribué
    à l'efficacité et à la précision de notre analyse.\\
    
    \noindent
    En termes de communication nous avons mis en place un espace Github relié à Github Desktop pour garder un environnement collaboratif. 
    Nous avons aussi pu échanger avec l’encadrant via des visioconférences Discord tout au long du projet.

\newpage
\section{Analyse des résultats}
    \emph{2 à 4 pages}
    \textcolor{blue}{
    L'analyse expérimentale sera faite dans cette section. Elle sera plus ou moins importante dans votre projet comme pour les deux points précédents, cette section devra donc être développée en conséquence.
    \begin{enumerate}% pour faire une liste numérotée
        \item Illustrer les performances ainsi que l'efficacité du logiciel implémenté à l'aide de graphiques.
        \item Analyser (et comparer, si plusieurs) les performances des solutions implémentées.
        \item Présenter les bancs d'essais (ou les procédures utilisées pour la génération des données) utilisés pour les tests du logiciel.
    \end{enumerate}
    }

\section{Gestion du Projet}
    \emph{(1-2 pages)}
    \textcolor{blue}{
    \begin{enumerate}% pour faire une liste numérotée
        \item Présenter la gestion du projet et les documents de planification rédigés (par exemple, le diagramme de Gantt).
        \item Discuter les changements majeurs effectués en cours de projet.
    \end{enumerate}
    }

\section{Bilan et Conclusions}
    \textcolor{blue}{
    Indiquer les fonctionnalités mises en œuvre par rapport au cahier des charges de départ, les points ouverts.
    Les perspectives du projet décrivent ce que l’on pourrait ajouter, modifier, réécrire afin d’améliorer le projet.
    }

\begin{thebibliography}{9}
    \bibitem{texbook}
    Donald E. Knuth (1986) \emph{The \TeX{} Book}, Addison-Wesley Professional.

    \bibitem{lamport94}
    Leslie Lamport (1994) \emph{\LaTeX: a document preparation system}, Addison
    Wesley, Massachusetts, 2nd ed.
    \end{thebibliography}

\end{document}
